\chapter{Discussion}\label{ch:discussion}
This chapter reflects on the reliability model introduced in \chref{ch:reliable-sots} and the demonstrative evaluation presented in \chref{ch:demonstration}. It first analyzes the utility of the approach (\secref{sec:discussion-utility}), then discusses implications for SoTS engineering practice (\secref{sec:discussion-implications}), and finally outlines future directions for extending the approach (\secref{sec:discussion-future}).


% \subsection{Non-Intrusive Integration with Existing Components}
% Statechart-based coordination can be introduced without requiring substantial modification to existing components. By acting as an external control layer, statecharts regulate participation and interaction without embedding coordination logic within services. This supports non-intrusive integration, allowing existing systems to be incorporated into a SoTS through lifecycle-aware wrappers. The approach aligns with ISO 23247, which emphasises interoperability and integration of heterogeneous systems without prescribing implementation details \cite{iso23247_2_2021}. It also supports multiple architectural viewpoints, enabling coordination to evolve independently of service implementation.


\section{Utility of Statechart-Based Coordination in SoTS}\label{sec:discussion-utility}
\subsection{Model-Driven Coordination and Behaviour Enforcement}

Results from \secref{sec:eval-statechart-enforcement} demonstrate that system behaviour is enforced directly through lifecycle statecharts, without requiring manual coordination logic. Lifecycle constraints define permissible configurations and transitions, ensuring that constituent behaviour remains consistent with system-level coordination policies at runtime.

This externalises coordination from an implicit, distributed responsibility within individual services to an explicit, shared model. As observed in the alignment between lifecycle states and system behaviour, constituents operate without embedding global coordination logic, yet system-level behaviour remains controlled and predictable. This establishes a clear separation of concerns between SoTS-level coordination and service-level functionality, which is particularly important in SoS contexts where operational and managerial independence limit tight integration \cite{maier1998architecting}.

By centralising coordination logic within executable models, the approach reduces implementation complexity and avoids duplication of coordination behaviour across components. Reliability is therefore achieved by construction through the coordination model, rather than through ad hoc safeguards within individual services. This aligns with model-driven engineering principles, where system behaviour is defined, enforced, and validated through models \cite{stahl2006model}. Statecharts further support the composition of multiple interacting lifecycles, enabling coordination to scale as system complexity increases \cite{harel1987statecharts}.


\subsection{Graceful Degradation and Adaptive Reliability}

As shown in \secref{sec:eval-graceful-degradation}, disturbances do not result in immediate system failure, but instead lead to progressive degradation of behaviour. As observation availability decreases, constituents transition between lifecycle states while maintaining partial functionality.

This behaviour corresponds to partial failure in dependability theory, where a system continues to deliver a subset of its intended service \cite{avizienis2004basic}. Lifecycle-driven transitions between participation modes (e.g., full, restricted, and passive roles) provide a structured mechanism for sustaining operation under degraded conditions.

Rather than reacting only after failure, the system adapts participation in response to changing conditions. Compensation mechanisms sustain event availability and maintain continuity of service without manual intervention, reflecting established resilience strategies in which degraded modes preserve essential functionality \cite{laprie1985dependable}. This shifts system behaviour from binary success/failure semantics toward controlled degradation, supporting key dependability attributes such as reliability and availability \cite{avizienis2004basic}.


\subsection{Statecharts as Executable Contracts for SoTS Coordination}

% Evidence from \secref{sec:eval-statechart-enforcement} and \secref{sec:eval-uncertainty-control} indicates that lifecycle constraints actively regulate behaviour at runtime. 
The use of statecharts to control constituent participation can be understood as establishing an explicit contract between the SoTS and each constituent. This contract defines how a constituent may contribute and how it is expected to behave under different conditions, making participation explicit rather than implicit or locally defined. Lifecycle states encode the assumptions under which participation is valid, while transitions define the guarantees regarding permissible actions and interactions. Constituents effectively adhere to these conditions through the externally defined lifecycle model, which regulates when and how they may contribute to system behaviour.

This interpretation aligns with contract-based design, where system composition is governed through assumptions and guarantees \cite{benveniste2012contracts}. However, unlike static interface specifications, these contracts are executable and enforced at runtime through state transitions that enable or restrict participation. As a result, lifecycle models provide both design-time and runtime assurances: they support analysis of properties such as reachability and consistency, while enforcing constraints during operation. System interaction therefore becomes explicit, enforceable, and amenable to compositional reasoning, with participation governed through operationalised contracts.




% \subsection{Statechart-Based Coordination as Executable Contracts}

% An important implication of the proposed approach is that lifecycle-based coordination establishes explicit contracts between the SoTS and its constituents. Rather than relying on implicit assumptions or locally implemented coordination logic, participation is controlled through a shared lifecycle model that defines how constituents may contribute under different conditions. Lifecycle states encode the assumptions under which participation is valid, while transitions define the guarantees regarding permissible actions and interactions, making coordination policies explicit and enforceable through externally defined behaviour.

% This perspective aligns with contract-based design, where system composition is governed through assumptions and guarantees \cite{benveniste2012contracts}. However, unlike traditional interface specifications, these contracts are executable and enforced at runtime through state transitions that enable or restrict participation. From an engineering perspective, this shifts coordination from informal agreements into a formally defined and enforceable mechanism, supporting both design-time analysis (e.g., reachability and consistency) and runtime assurance. As a result, coordination becomes explicit, verifiable, and amenable to compositional reasoning, improving the reliability and governability of SoTS interactions.



\section{Implications for SoTS Engineering Practice}\label{sec:discussion-implications}
The proposed approach has several implications for SoTS engineering, particularly in how coordination, information exchange, and system-level interaction are structured across independently developed components.

\subsection{Privacy-Preserving Coordination through Minimal Information Exposure}
A key implication is the separation of concerns between component internals and system-level interaction. Constituents participate through constrained lifecycle states without exposing internal logic, enabling privacy-preserving coordination aligned with GDPR principles such as data minimisation and purpose limitation \cite{voigt2017eu}. Because participation is governed by lifecycle constraints, only information required to determine eligibility and progression is exposed, limiting interaction to minimal, purpose-specific exchange and reducing unnecessary disclosure by design. This perspective is reflected in SoTS research such as \citet{heininger2021capturing}, which embeds privacy management in digital twin models through controlled data sharing and selective disclosure.

An additional consequence is reduced cognitive burden for developers. For example, a CEP engineer interacts only with a defined set of lifecycle states and events, without reasoning about the full SoTS lifecycle or internal behaviour of other constituents. This allows domain specialists to focus on local functionality while the coordination model mediates both interaction and information exposure, becoming the primary mechanism governing system behaviour and data visibility.

\subsection{Lifecycle-Based Control as an Enabler of Engineering Maturity}

Lifecycle-based control enables more mature and systematic engineering practices by making system-level interaction explicit and structured. Coordination concerns that are often implicit or ad hoc are externalised into lifecycle models, allowing analysis, validation, and governance. This supports different maturity levels: in less mature contexts, lifecycle models introduce standardisation through a shared representation of participation and interaction, while in more advanced settings they integrate into toolchains for code generation, validation, and runtime monitoring, enabling more automated and controlled system evolution. This progression aligns with maturity models such as CMM \cite{paulk1991capability} and reflects the importance of organisational capability in SoS engineering \cite{carlock2001system}.

Adopting this approach requires organisational capabilities in modelling, state-based reasoning, and system-level design, along with governance structures to manage evolving lifecycle definitions and coordination policies. It may also introduce new roles, such as SoTS architects or model engineers. Beyond their technical role, executable lifecycle models act as shared reference points that support communication across stakeholders and enable more systematic and controlled engineering processes.

\section{Future Directions for State-Based SoTS Coordination}\label{sec:discussion-future}

The proposed approach provides a foundation for state-based coordination in SoTS, but several extensions can further improve its expressiveness, adaptability, and analytical capabilities.

\subsection{Extending Lifecycle-Based Coordination}

Future work could extend the lifecycle-based model by incorporating additional orthogonal dimensions that influence participation. Trust relationships may determine whether a constituent is eligible to participate within certain states, particularly in dynamic or federated environments. Operational readiness can capture whether a system is prepared to contribute based on configuration or resource availability, while data integrity ensures that only systems providing valid and consistent information are involved in ongoing activities. Additional dimensions such as security and quality-of-service (QoS) \cite{o2007quality} can further refine participation semantics by restricting involvement based on authorization and performance requirements.

These extensions point toward multi-dimensional state models, where orthogonal concerns refine lifecycle semantics and constrain participation in a structured manner. Embedding these dimensions directly into the state representation strengthens the abstraction of SoTS structure and behaviour, supporting more precise and context-aware control over which constituents are active, trusted, and effective in each lifecycle phase.

\subsection{Dynamic Selection of Fault Tolerance Strategies}

The current approach assumes a fixed fault tolerance or compensation strategy, which may not be optimal under varying runtime conditions. Future systems could instead support dynamic selection of fault tolerance mechanisms based on system context, including resource availability, failure severity, and performance requirements. This would allow coordination to adapt more effectively to changing environments. A coordination framework could maintain a portfolio of strategies---such as retry mechanisms, compensation actions, or redundancy---and select among them at runtime. This selection could be guided by policies or utility functions that balance dependability attributes \cite{avizienis2004basic}, resulting in more flexible fault handling. Learning-based approaches, including reinforcement learning, could further improve decision-making over time, enabling fault tolerance to evolve based on observed system behaviour \cite{dong2023deep, zhang2020fault}.

\subsection{Resource-Aware Coordination through Petri Nets}

An important extension of the proposed approach is the integration of resource-aware modelling into lifecycle-based representations. While statecharts can encode resource availability and capacity constraints through variables and guards, these aspects are often implicit and embedded within control logic, making explicit reasoning about resource usage and contention more difficult in complex SoTS environments. Petri nets provide a complementary formalism for modelling resources, concurrency, and synchronization \cite{peterson1981petri}. By representing resources as tokens, they enable direct reasoning about availability and consumption while naturally capturing parallelism and shared resource access. This improves transparency and analyzability, and has been applied in SoTS contexts to model process flows and interactions \cite{zhang2022multi-scale}.

Petri nets have also been used to provide formal semantics for state-based models, supporting enhanced simulation and analysis of dynamic behaviour \cite{baresi2001improving}. Combining statecharts with Petri net models enables coordination decisions to be grounded in both lifecycle state and resource conditions. Transitions can be guarded by resource constraints derived from the Petri net model, allowing participation and progression to depend on both resource availability and behavioural state. This hybrid approach results in a more expressive and analyzable model of SoTS coordination.


\subsection{Verification of Statechart-Based SoTS Coordination}

As SoTS are increasingly deployed in complex and dynamic environments, ensuring the correctness of coordination mechanisms becomes essential. State-based models provide a strong foundation for this, as they explicitly define permissible system behaviour through states and transitions. This enables verification of key properties, such as preventing unreliable contributions and avoiding unsafe participation states. These properties can be systematically analyzed using automated safety analysis techniques \cite{pap2005methods} and formal verification methods such as reachability analysis \cite{kang1996formalization}. The model also supports traceability between coordination rules and observed behaviour, allowing verification results to be mapped to lifecycle states and their associated coordination rules. Future work could extend this by applying formal verification to prove safety and liveness properties, as well as integrating runtime monitoring to ensure coordination constraints remain enforced during execution in dynamic SoTS environments.